\documentclass[11pt,a4paper]{article}
\usepackage[margin=2cm]{geometry}
\usepackage[T1]{fontenc}
\usepackage{lmodern}
\usepackage{booktabs,longtable,array}
\usepackage{xcolor}
\usepackage{hyperref}
\hypersetup{colorlinks=true,linkcolor=blue,urlcolor=blue}
\title{OpsDesk\\Real-Time Service Request Management System}
\author{}
\date{}
\begin{document}
\maketitle

\section{Overview}
OpsDesk is a full-stack service request management system. Operators submit long-running tasks without waiting for completion. Supervisors monitor status, progress, queue activity, and audit logs in real time through WebSockets and background Worker Threads.

\section{Features}
Non-blocking ingestion; bounded multi-threaded processing; Socket.IO live updates; Operator and Supervisor workspaces; cooperative cancellation; MongoDB progress audit trails; late-join hydration; search, filters, pagination; and table/grid Supervisor views.

\section{Technology Stack}
\begin{longtable}{p{0.25\textwidth}p{0.65\textwidth}}
\toprule
Frontend & React 19, TypeScript, Vite, Tailwind CSS, TanStack React Query, Zustand, Lucide, Socket.IO Client\\
Backend & Node.js, Express, TypeScript, Socket.IO, Zod, CORS, express-rate-limit\\
Concurrency & Node.js \texttt{worker\_threads}\\
Database & MongoDB with Mongoose\\
Testing & Vitest, Supertest, MongoMemoryServer\\
\bottomrule
\end{longtable}

\section{Prerequisites}
Install Node.js 18 or later, npm 9 or later, and MongoDB Community Server or use MongoDB Atlas. The default database URI is \texttt{mongodb://localhost:27017/service-requests}.

\section{Setup and Environment}
Create \texttt{server/.env}:
\begin{verbatim}
NODE_ENV=development
PORT=5000
MONGODB_URI=mongodb://localhost:27017/service-requests
CORS_ORIGIN=http://localhost:5173
MAX_WORKERS=5
RATE_LIMIT_WINDOW_MS=60000
RATE_LIMIT_MAX_REQUESTS=20
\end{verbatim}
Install dependencies:
\begin{verbatim}
cd server
npm install
cd ../client
npm install
\end{verbatim}

\section{Build and Run}
Start MongoDB, then use two terminals:
\begin{verbatim}
# Terminal 1
cd server
npm run dev

# Terminal 2
cd client
npm run dev
\end{verbatim}
Open \url{http://localhost:5173}. The backend runs at \url{http://localhost:5000}. If port 5000 is already in use, keep the healthy existing backend or change \texttt{PORT} and the Vite proxy together.

Production commands:
\begin{verbatim}
cd server && npm run build && npm start
cd client && npm run build && npm run preview
\end{verbatim}

\section{Database Setup}
The server connects to MongoDB during startup. It stores service requests in the \texttt{ServiceRequest} collection and checkpoint records in the \texttt{ProgressLog} collection. Indexes support status, priority, submitter, timestamps, and text search. Tests use an isolated in-memory MongoDB instance.

\section{API Documentation}
\begin{longtable}{p{0.12\textwidth}p{0.35\textwidth}p{0.38\textwidth}}
\toprule Method & Endpoint & Description\\\midrule
POST & \texttt{/api/requests} & Create and enqueue a service request\\
GET & \texttt{/api/requests} & List, search, filter, sort, and paginate requests\\
GET & \texttt{/api/requests/:id} & Retrieve request details\\
POST & \texttt{/api/requests/:id/cancel} & Cancel pending or processing request\\
GET & \texttt{/api/requests/:id/progress} & Retrieve progress audit logs\\
GET & \texttt{/health} & Return server health\\
\bottomrule
\end{longtable}
Creation requires \texttt{title}, \texttt{description}, \texttt{priority}, and \texttt{submittedBy}. Valid priorities are \texttt{low}, \texttt{medium}, \texttt{high}, and \texttt{critical}. List queries support \texttt{page}, \texttt{limit}, \texttt{status}, \texttt{priority}, \texttt{submittedBy}, \texttt{search}, \texttt{sortBy}, and \texttt{sortOrder}.

\section{WebSocket Events}
Socket.IO broadcasts \texttt{request:created}, \texttt{request:status-updated}, \texttt{request:progress-updated}, \texttt{request:completed}, \texttt{request:failed}, and \texttt{request:cancelled}. A newly connected client receives \texttt{requests:initial-state} containing pending and processing requests.

\section{Testing}
Run the backend API tests with:
\begin{verbatim}
cd server
npm test
\end{verbatim}
The tests cover validation, persistence, listing, filtering, pagination, cancellation, progress-log retrieval, and errors. Worker Threads and Socket.IO are mocked in the API suite; live worker and WebSocket behavior should also be demonstrated manually in two browser views.

\section{Assumptions and Limitations}
Roles are simulated with a client-side Operator/Supervisor toggle; authentication and backend RBAC are not implemented. The WorkerPool queue is in memory and the deployment is single-node. Authentication, distributed queuing, and external notification webhooks are future improvements.
\end{document}
